\section{Gobernanza del producto}
La gobernanza del ecosistema documental define responsabilidades, criterios de aprobación y mecanismos de control para asegurar continuidad institucional y coherencia evolutiva.

\section{Modelo de decision}
\begin{itemize}
  \item Comité funcional: prioriza requerimientos de Archivo y RRHH.
  \item Comité técnico: valida arquitectura, seguridad y viabilidad.
  \item Coordinación institucional: aprueba alcance, cronograma y riesgos.
\end{itemize}

\section{Politicas de versionado y cambio}
\begin{itemize}
  \item Versionado semántico de entregables técnicos y documentales.
  \item Control de cambios con evidencia de impacto.
  \item Registro de aprobaciones y decisiones arquitectónicas.
\end{itemize}

\section{Marco legal y cumplimiento}
\begin{itemize}
  \item Cumplimiento de lineamientos institucionales de privacidad.
  \item Resguardo de información sensible y de carácter personal.
  \item Trazabilidad de operaciones para auditoría interna.
  \item Politicas de propiedad intelectual y uso autorizado.
\end{itemize}

\section{RACI de alto nivel}
\begin{longtable}{p{0.28\textwidth} p{0.16\textwidth} p{0.16\textwidth} p{0.16\textwidth} p{0.16\textwidth}}
\toprule
\textbf{Proceso} & \textbf{R} & \textbf{A} & \textbf{C} & \textbf{I} \\
\midrule
Priorización de backlog & Equipo funcional & Coordinación proyecto & Equipo técnico & Autoridades \\
Aprobación de despliegue & Equipo técnico & Coordinación proyecto & Seguridad/QA & Usuarios clave \\
Actualizacion documental & Equipo documental & Coordinacion proyecto & Equipo funcional & Autoridades \\
Auditoría y trazabilidad & Seguridad/Control & Coordinación institucional & Equipo técnico & Autoridades \\
\bottomrule
\end{longtable}
